\section{Đảm bảo chất lượng dữ liệu}
\label{sec:data-quality}

\subsection{Tổng quan chất lượng}

Sau quá trình thu thập và làm sạch, chất lượng dữ liệu được đánh giá theo bốn tiêu chí: tính đầy đủ (completeness), tính nhất quán (consistency), tính chính xác (accuracy), và tính kịp thời (timeliness).

\begin{table}[H]
    \centering
    \caption{Đánh giá chất lượng dữ liệu theo nguồn}
    \label{tab:data-quality}
    \begin{tabular}{lcc}
        \toprule
        \textbf{Tiêu chí} & \textbf{Chợ Tốt} & \textbf{Websosanh} \\
        \midrule
        Completeness (trung bình) & $\sim$85\% & $\sim$93\% \\
        Consistency (brand--title) & $\sim$95\% & $\sim$98\% \\
        Accuracy (giá hợp lệ) & $>$99\% & $>$95\% (sau lọc outlier) \\
        Timeliness & Q1 2025 & Q1 2025 \\
        \bottomrule
    \end{tabular}
\end{table}

\subsection{Xử lý outlier}

\subsubsection{Outlier giá --- Chợ Tốt}
Giá đăng bán có range từ 100.000 VND đến trên 150 triệu VND. Các giá cực thấp ($<$ 1 triệu) thường là phụ kiện hoặc linh kiện bị gắn nhầm category. Giá cực cao ($>$ 100 triệu) phần lớn hợp lệ (Apple MacBook Pro cấu hình cao). Nhóm nghiên cứu áp dụng domain filter: giữ lại giá trong khoảng 1--150 triệu VND.

\subsubsection{Outlier giá --- Websosanh}
Dữ liệu chứa các giá bất thường (ví dụ 268 triệu, 860 triệu, 1.629 triệu, 2.750 triệu VND) do lỗi crawl hoặc parse. Quy trình lọc gồm hai bước:
\begin{enumerate}[noitemsep]
    \item \textbf{Domain filter}: Null-out giá $<$ 3 triệu hoặc $>$ 200 triệu.
    \item \textbf{Relative filter}: Dùng median giá của từng dòng làm anchor; null-out giá lệch $>$ 2$\times$ hoặc $<$ 0.5$\times$ so với row median.
\end{enumerate}

Sau cleaning, mỗi sản phẩm Websosanh trung bình còn $\sim$2.99 nguồn giá sạch.

\subsubsection{Outlier thuộc tính}
\begin{itemize}[noitemsep]
    \item \textbf{Kích thước màn hình}: Giá trị $>$ 25 inch được null-out (không phù hợp laptop).
    \item \textbf{RAM}: Giá trị 512 hoặc 5.600 GB là lỗi parse (nhầm tốc độ RAM với dung lượng); null-out giá trị $>$ 256 GB.
    \item \textbf{Storage}: Giá trị 6.144 GB ($\approx$ 6TB) kiểm tra thủ công --- một số là cấu hình thật (cộng gộp nhiều ổ).
\end{itemize}

\subsection{Tạo biến phái sinh (Feature Engineering)}

Để giảm phân mảnh (fragmentation) và sparse features, các biến categorical được gom nhóm:

\begin{table}[H]
    \centering
    \caption{Các biến phái sinh chính}
    \label{tab:derived-features}
    \begin{tabular}{lll}
        \toprule
        \textbf{Biến gốc} & \textbf{Biến mới} & \textbf{Nhóm} \\
        \midrule
        Kích thước (inch)  & \texttt{screen\_size\_group} & $<$13, 13--13.9, 14--14.9, 15--15.9, 16--16.9, $\geq$17 \\
        Dung lượng RAM     & \texttt{ram\_tier\_clean}    & $\leq$8GB, 16GB, 32GB, $>$32GB \\
        Dung lượng ổ cứng  & \texttt{storage\_tier\_clean} & $\leq$256GB, 512GB, 1TB, 2TB, $>$2TB \\
        GPU model          & \texttt{gpu\_tier\_v2}       & Intel Integrated, AMD Radeon, GTX, RTX 4000, RTX 5000, Apple, Other \\
        Loại CPU           & \texttt{cpu\_tier}           & Low-end, Mid-range, High-end, Unknown \\
        \bottomrule
    \end{tabular}
\end{table}

\subsection{Chất lượng biến Unknown}

Sau toàn bộ quá trình cleaning, tỷ lệ Unknown còn lại ở Websosanh:
\begin{itemize}[noitemsep]
    \item \texttt{ram\_type\_clean}: 7.44\%
    \item \texttt{cpu\_tier}: 6.93\% (304 dòng)
    \item \texttt{storage\_tier\_clean}: 5.13\%
    \item \texttt{storage\_type\_clean}: 5.06\%
    \item \texttt{gpu\_tier\_v2}: 4.72\% (207 dòng)
    \item \texttt{brand\_clean}: 4.11\%
\end{itemize}

Các nhóm Unknown được giữ lại như một category riêng thay vì xóa dòng, nhằm tránh mất mẫu và bias.
